%%%%%%%%%%%%%%%%%%%%%%%%%%%%%%%%%%%%%%%%%%%%%%%%%%%%%%%%%%%%%%%%%%%%%%
% Lớp 11 | Quan hệ vuông góc trong không gian | Hình lăng trụ đứng. Hình chóp đều. Thể tích
% Dạng 1H8?7-7: Tỉ số thể tích
% Số câu: 3  (sinh tự động bởi tach_cau.py; từ nay thầy sửa trực tiếp file này)
%%%%%%%%%%%%%%%%%%%%%%%%%%%%%%%%%%%%%%%%%%%%%%%%%%%%%%%%%%%%%%%%%%%%%%

%%%%%%%%%%%%%%%% TRẮC NGHIỆM 4 PHƯƠNG ÁN %%%%%%%%%%%%%%%%

% nguồn: Ok_Toan11_TracNghiem_Lan2
\begin{ex}%[1H8N7-7]%[ma:f3cff4bb]
	Cho khối chóp và khối lăng trụ có diện tích đáy, chiều cao tương ứng  bằng nhau và có thể tích lần lượt là $V_1$, $V_2$. Tỉ số $\dfrac{V_1}{V_2}$ bằng
	\choice 
		 {$\dfrac{2}{3}$}
		{$\dfrac{3}{2}$} 
		{$3$} 
		{\True  $\dfrac{1}{3}$}
	\loigiai{ 
		Gọi đường cao, diện tích đáy  lần lượt là $h,B$.\\
		Khi đó áp dụng công thức thể tích khối chóp, khối lăng trụ  ta được $V_1=\dfrac{1}{3}B\cdot h$ và $V_2=B\cdot h$.\\
	Suy ra $\dfrac{V_1}{V_2}=\dfrac{\dfrac{1}{3}B\cdot h}{B\cdot  h}=\dfrac{1}{3}$.}
\end{ex}

% nguồn: Ok_Toan11_TracNghiem_Lan2
\begin{ex}%[1H8V7-7]%[ma:bb94d46d]
	Cho tứ diện $ABCD$ có các cạnh $AB$, $AC$ và $AD$ đôi một vuông góc với nhau; $AB=6a$, $AC=7a$ và $AD=4a$. Gọi $M$, $N$, $P$ tương ứng là trung điểm các cạnh $BC$, $CD$, $DB$. Tính thể tích  $V$ của tứ diện $AMNP$.
	 \choice
		 {\True $V=7a^3$}
		{$V=14a^3$} 
		 {$V=\dfrac{28}{3}a^3$}
		{$V=\dfrac{7}{2}a^3$} 
	\loigiai{ 
		\begin{center}
			\begin{tikzpicture}[scale=1,font=\footnotesize,line join=round,line cap=round,>=stealth]
					\coordinate[label=left:$A$] (A)  at (0,0);
					\coordinate[label=right:$B$]  (B) at (4,0);
					\coordinate[label=below left:$C$] (C) at  (-40:4);
					\coordinate[label=above:$D$] (D) at  (90:4);
					\coordinate[label=right:$M$]  (M) at ($(C)!0.5!(B)$);
					\coordinate[label=right:$N$]  (N) at ($(C)!0.5!(D)$);
					\coordinate[label=right:$P$] (P) at  ($(D)!0.5!(B)$);
					\draw (A)--(C)--(B)--(D)--cycle  (D)--(C)
						 (A)--(N)--(M)--(P)--(N);
					\draw[dashed] (A)--(B) 
						 (P)--(A)--(M);
					\foreach  \diem in {A,C,B,D,M,N,P} \fill (\diem) circle (1pt);
					 \draw pic[draw,angle radius=0.3cm]{right angle=M--A--D};
					\draw  pic[draw,angle radius=0.3cm]{right angle=B--M--A};
				\end{tikzpicture}
		\end{center}
		Vì $AB$, $AC$, $AD$ đôi một vuông góc, thể tích của  tứ diện $ABCD$ bằng
		$$V_{ABCD} = \dfrac{1}{6} \cdot AB \cdot AC  \cdot AD = \dfrac{1}{6} \cdot 6a \cdot 7a \cdot 4a = 28a^3.$$
		 Ta có
		$S_{MNP} = \dfrac{1}{2}\cdot S_{MNPD} = \dfrac{1}{4}\cdot S_{BCD}$.\\ 
		Suy ra 
		 $$V_{AMNP} = \dfrac{1}{4}\cdot V_{ABCD}.$$
		Khi đó, ta  được
		$$V_{AMNP} = \dfrac{1}{4}\cdot V_{ABCD} = \dfrac{1}{4}\cdot  28a^3 = 7a^3.$$
	 }
\end{ex}

% nguồn: Ok_Toan11_TracNghiem_Lan2
\begin{ex}%[1H8C7-7]%[ma:8bdd6b54]
	Cho khối tứ diện có  thể tích bằng $V$. Gọi $V'$ là thể tích của khối đa diện có các đỉnh là các trung điểm của các cạnh của khối tứ diện đã cho, tính tỉ số $\dfrac{V'}{V}$.
	 \choice
		 {\True $\dfrac{V'}{V} = \dfrac{1}{2}$}
		{$\dfrac{V'}{V} = \dfrac{1}{4}$} 
		{$\dfrac{V'}{V}  = \dfrac{2}{3}$}
		{$\dfrac{V'}{V} = \dfrac{5}{8}$} 
	\loigiai{ 
		\begin{center}
			\begin{tikzpicture}[scale=1,font=\footnotesize,line join=round,line cap=round,>=stealth]
					\coordinate[label=left:$A$] (A) at  (0,0);
					 \coordinate[label=right:$B$] (B) at (4,0);
					\coordinate[label=below left:$C$] (C) at (-40:4); 
					\coordinate[label=above:$D$] (D) at  (50:4);
					\coordinate[label=right:$M$] (M)  at ($(C)!0.5!(B)$);
					\coordinate[label=right:$N$]  (N) at ($(C)!0.5!(D)$);
					 \coordinate[label=right:$P$] (P) at ($(D)!0.5!(B)$);
					 \coordinate[label=below:$Q$] (Q) at ($(A)!0.5!(B)$);
					\coordinate[label=below:$R$]  (R) at ($(A)!0.5!(C)$);
					\coordinate[label=left:$T$] (T) at  ($(A)!0.5!(D)$);
					\draw (A)--(C)--(B)--(D)--cycle (D)--(C) 
						 (M)--(P)--(N)--cycle (R)--(N)--(T)--cycle;
					 \draw[dashed] (A)--(B)
						(Q)--(T)--(P)--(Q)--(M)--(R)--cycle; 
					\foreach \diem in {A,C,B,D,M,N,P,Q,R,T}  \fill (\diem) circle (1pt);
				\end{tikzpicture}
		\end{center}
		 Khối đa diện có các đỉnh là trung điểm các cạnh của một tứ diện là một khối tám mặt đều.\\
		Thể tích $V'$ của khối  tám mặt này có thể được tính bằng cách lấy thể tích $V$ của khối tứ diện ban đầu trừ đi thể tích của $4$ khối tứ diện nhỏ ở $4$ đỉnh.\\
		Xét một đỉnh bất kỳ, ví dụ đỉnh $A$ của tứ diện $ABCD$. Khối tứ diện nhỏ  ở đỉnh $A$ có các đỉnh là $A$ và trung điểm của các cạnh $AB$, $AC$, $AD$.
		Thể tích của  khối tứ diện nhỏ này bằng
		$$V_{\text{nhỏ}} = \dfrac{1}{8}\cdot  V.$$
		Có 4 khối tứ diện nhỏ ở 4 đỉnh, do đó tổng thể tích của chúng là 
		$$4  \cdot \dfrac{1}{8}\cdot V = \dfrac{1}{2}V.$$
		 Thể tích khối tám mặt ở giữa là
		$$V' = V - 4 \cdot V_{\text{nhỏ}} = V - \dfrac{1}{2}V =  \dfrac{1}{2}V.$$
		Vậy tỉ  số cần tìm là $\dfrac{V'}{V} = \dfrac{1}{2}$.
	 }
\end{ex}
